% =====================================================================
%  Chapitre « Présentation générale d'OpenStack »
%  Fichier importé par main.tex avec % =====================================================================
%  Chapitre « Présentation générale d'OpenStack »
%  Fichier importé par main.tex avec % =====================================================================
%  Chapitre « Présentation générale d'OpenStack »
%  Fichier importé par main.tex avec % =====================================================================
%  Chapitre « Présentation générale d'OpenStack »
%  Fichier importé par main.tex avec \include{chap_presentation}.
%  Il ne contient ni préambule ni \begin{document} : les environnements
%  (definition, resultat, remarque, aretenir, savoirfaire, exercice, correction),
%  les styles TikZ (svc, svcgris, bus, flux) et les couleurs (insarouge,
%  insagris) sont définis dans main.tex. Un chapitre = une \section.
%  Contenu à rédiger.
% =====================================================================

\section{Présentation générale d'OpenStack}\label{chap:presentation}

% Squelette tiré du support « Introduction à OpenStack » : titres de sous-sections seulement.

% ---------------------------------------------------------------------
\subsection{Du serveur physique au cloud}\label{sec:pres-cloud}
% ---------------------------------------------------------------------
% virtualisation, machines virtuelles, hyperviseur

% ---------------------------------------------------------------------
\subsection{Modèles de service et de déploiement}\label{sec:pres-modeles}
% ---------------------------------------------------------------------
% IaaS, PaaS, SaaS ; cloud privé, public, hybride

% ---------------------------------------------------------------------
\subsection{Qu'est-ce qu'OpenStack ?}\label{sec:pres-openstack}
% ---------------------------------------------------------------------
% plateforme IaaS open source ; OpenStack n'est pas un hyperviseur

% ---------------------------------------------------------------------
\subsection{Architecture générale}\label{sec:pres-architecture}
% ---------------------------------------------------------------------
% vue d'ensemble des services, API REST, bus de messages, microservices

% ---------------------------------------------------------------------
\subsection{Architecture multi-nœuds}\label{sec:pres-multinoeuds}
% ---------------------------------------------------------------------
% rôles des nœuds, séparation du trafic réseau

% ---------------------------------------------------------------------
\subsection{Que se passe-t-il quand on crée une VM ?}\label{sec:pres-scenario}
% ---------------------------------------------------------------------
% scénario complet, de la commande CLI aux services sollicités

% ---------------------------------------------------------------------
\subsection{OpenStack et les technologies voisines}\label{sec:pres-comparaison}
% ---------------------------------------------------------------------
% KVM, Proxmox, VMware vSphere, Kubernetes
.
%  Il ne contient ni préambule ni \begin{document} : les environnements
%  (definition, resultat, remarque, aretenir, savoirfaire, exercice, correction),
%  les styles TikZ (svc, svcgris, bus, flux) et les couleurs (insarouge,
%  insagris) sont définis dans main.tex. Un chapitre = une \section.
%  Contenu à rédiger.
% =====================================================================

\section{Présentation générale d'OpenStack}\label{chap:presentation}

% Squelette tiré du support « Introduction à OpenStack » : titres de sous-sections seulement.

% ---------------------------------------------------------------------
\subsection{Du serveur physique au cloud}\label{sec:pres-cloud}
% ---------------------------------------------------------------------
% virtualisation, machines virtuelles, hyperviseur

% ---------------------------------------------------------------------
\subsection{Modèles de service et de déploiement}\label{sec:pres-modeles}
% ---------------------------------------------------------------------
% IaaS, PaaS, SaaS ; cloud privé, public, hybride

% ---------------------------------------------------------------------
\subsection{Qu'est-ce qu'OpenStack ?}\label{sec:pres-openstack}
% ---------------------------------------------------------------------
% plateforme IaaS open source ; OpenStack n'est pas un hyperviseur

% ---------------------------------------------------------------------
\subsection{Architecture générale}\label{sec:pres-architecture}
% ---------------------------------------------------------------------
% vue d'ensemble des services, API REST, bus de messages, microservices

% ---------------------------------------------------------------------
\subsection{Architecture multi-nœuds}\label{sec:pres-multinoeuds}
% ---------------------------------------------------------------------
% rôles des nœuds, séparation du trafic réseau

% ---------------------------------------------------------------------
\subsection{Que se passe-t-il quand on crée une VM ?}\label{sec:pres-scenario}
% ---------------------------------------------------------------------
% scénario complet, de la commande CLI aux services sollicités

% ---------------------------------------------------------------------
\subsection{OpenStack et les technologies voisines}\label{sec:pres-comparaison}
% ---------------------------------------------------------------------
% KVM, Proxmox, VMware vSphere, Kubernetes
.
%  Il ne contient ni préambule ni \begin{document} : les environnements
%  (definition, resultat, remarque, aretenir, savoirfaire, exercice, correction),
%  les styles TikZ (svc, svcgris, bus, flux) et les couleurs (insarouge,
%  insagris) sont définis dans main.tex. Un chapitre = une \section.
%  Contenu à rédiger.
% =====================================================================

\section{Présentation générale d'OpenStack}\label{chap:presentation}

% Squelette tiré du support « Introduction à OpenStack » : titres de sous-sections seulement.

% ---------------------------------------------------------------------
\subsection{Du serveur physique au cloud}\label{sec:pres-cloud}
% ---------------------------------------------------------------------
% virtualisation, machines virtuelles, hyperviseur

% ---------------------------------------------------------------------
\subsection{Modèles de service et de déploiement}\label{sec:pres-modeles}
% ---------------------------------------------------------------------
% IaaS, PaaS, SaaS ; cloud privé, public, hybride

% ---------------------------------------------------------------------
\subsection{Qu'est-ce qu'OpenStack ?}\label{sec:pres-openstack}
% ---------------------------------------------------------------------
% plateforme IaaS open source ; OpenStack n'est pas un hyperviseur

% ---------------------------------------------------------------------
\subsection{Architecture générale}\label{sec:pres-architecture}
% ---------------------------------------------------------------------
% vue d'ensemble des services, API REST, bus de messages, microservices

% ---------------------------------------------------------------------
\subsection{Architecture multi-nœuds}\label{sec:pres-multinoeuds}
% ---------------------------------------------------------------------
% rôles des nœuds, séparation du trafic réseau

% ---------------------------------------------------------------------
\subsection{Que se passe-t-il quand on crée une VM ?}\label{sec:pres-scenario}
% ---------------------------------------------------------------------
% scénario complet, de la commande CLI aux services sollicités

% ---------------------------------------------------------------------
\subsection{OpenStack et les technologies voisines}\label{sec:pres-comparaison}
% ---------------------------------------------------------------------
% KVM, Proxmox, VMware vSphere, Kubernetes
.
%  Il ne contient ni préambule ni \begin{document} : les environnements
%  (definition, resultat, remarque, aretenir, savoirfaire, exercice, correction),
%  les styles TikZ (svc, svcgris, bus, flux) et les couleurs (insarouge,
%  insagris) sont définis dans main.tex. Un chapitre = une \section.
%  Contenu à rédiger.
% =====================================================================

\section{Présentation générale d'OpenStack}\label{chap:presentation}

% Squelette tiré du support « Introduction à OpenStack » : titres de sous-sections seulement.

% ---------------------------------------------------------------------
\subsection{Du serveur physique au cloud}\label{sec:pres-cloud}
% ---------------------------------------------------------------------
% virtualisation, machines virtuelles, hyperviseur

% ---------------------------------------------------------------------
\subsection{Modèles de service et de déploiement}\label{sec:pres-modeles}
% ---------------------------------------------------------------------
% IaaS, PaaS, SaaS ; cloud privé, public, hybride

% ---------------------------------------------------------------------
\subsection{Qu'est-ce qu'OpenStack ?}\label{sec:pres-openstack}
% ---------------------------------------------------------------------
% plateforme IaaS open source ; OpenStack n'est pas un hyperviseur

% ---------------------------------------------------------------------
\subsection{Architecture générale}\label{sec:pres-architecture}
% ---------------------------------------------------------------------
% vue d'ensemble des services, API REST, bus de messages, microservices

% ---------------------------------------------------------------------
\subsection{Architecture multi-nœuds}\label{sec:pres-multinoeuds}
% ---------------------------------------------------------------------
% rôles des nœuds, séparation du trafic réseau

% ---------------------------------------------------------------------
\subsection{Que se passe-t-il quand on crée une VM ?}\label{sec:pres-scenario}
% ---------------------------------------------------------------------
% scénario complet, de la commande CLI aux services sollicités

% ---------------------------------------------------------------------
\subsection{OpenStack et les technologies voisines}\label{sec:pres-comparaison}
% ---------------------------------------------------------------------
% KVM, Proxmox, VMware vSphere, Kubernetes
